\clearpage\pagestyle{otMain}\setlength{\otBlocEspacement}{2.5mm}\fontsize{11}{15}\selectfont\corrigesChapitre{9}{Comparer et arrondir les décimaux}
\begin{blocLivre}{cor-9-1}\begin{corrige}{1}{Le piège classique}\demarche{Écrire les deux nombres avec deux chiffres après la virgule.}
$3{,}5=3{,}50$ : cinquante centièmes contre quarante-cinq. Donc $\res{3{,}5>3{,}45}$.
\end{corrige}\end{blocLivre}
\begin{blocLivre}{cor-9-2}\begin{corrige}{2}{Compléter par un symbole}\demarche{Comparer rang par rang depuis la gauche.}
Parties entières égales puis $0<8$ aux dixièmes : $\res{6{,}08<6{,}8}$. Les parties entières décident ensuite : $\res{12{,}34>9{,}99}$.
\end{corrige}\end{blocLivre}
\begin{blocLivre}{cor-9-3}\begin{corrige}{3}{Ranger cinq nombres}\demarche{Aligner mentalement les rangs décimaux.}
$6{,}99$ a la plus petite partie entière. Ensuite, comparer les dixièmes puis les centièmes : $\res{6{,}99<7{,}09<7{,}1<7{,}3<7{,}31}$.
\end{corrige}\end{blocLivre}
\begin{blocLivre}{cor-9-4}\begin{corrige}{4}{Encadrer par des entiers}\demarche{Lire la partie entière.}
$\res{4<4{,}62<5}$ et $\res{0<0{,}85<1}$. La partie entière donne la borne de gauche.
\end{corrige}\end{blocLivre}
\begin{blocLivre}{cor-9-5}\begin{corrige}{5}{Partie entière}\demarche{Prendre l’entier juste avant le nombre.}
Les parties entières sont $\res{17}$, $\res9$ et $\res0$.
\end{corrige}\end{blocLivre}
\begin{blocLivre}{cor-9-6}\begin{corrige}{6}{Lire entre les graduations}\demarche{Identifier la valeur d’un pas avant de lire.}
$\res{A(5{,}3)}$ et $\res{B(5{,}8)}$. Le point C se place au milieu du pas entre $5{,}4$ et $5{,}5$ ; il figure en vert sur le corrigé.
\label{fig-cor-9-6}\figureGeo[corrige]{dec9-6}
\end{corrige}\end{blocLivre}
\begin{blocLivre}{cor-9-7}\begin{corrige}{7}{Intercaler}\demarche{Ajouter un rang décimal pour affiner.}
En centièmes : $\res{2{,}45}$ convient, car $2{,}40<2{,}45<2{,}50$. En ajoutant un rang : $\res{8{,}315}$ convient ; tout nombre strictement entre les bornes est accepté.
\end{corrige}\end{blocLivre}
\begin{blocLivre}{cor-9-8}\begin{corrige}{8}{Arrondir à l’entier}\demarche{Regarder le chiffre des dixièmes.}
$7{,}3$ est plus proche de 7 : arrondi $\res7$. $7{,}83$ est plus proche de 8 : arrondi $\res8$. $12{,}5$ est à égale distance : on choisit $\res{13}$.
\end{corrige}\end{blocLivre}
\begin{blocLivre}{cor-9-9}\begin{corrige}{9}{Compter de 0,1 en 0,1}\demarche{Échanger dix dixièmes contre une unité.}
$3{,}7$ ; $3{,}8$ ; $3{,}9$ ; $\res{4}$ ; $4{,}1$. Dix dixièmes font une unité : après $3{,}9$ vient 4, et non « $3{,}10$ ».
\end{corrige}\end{blocLivre}
\begin{blocLivre}{cor-9-10}\begin{corrige}{10}{Le même nombre ?}\demarche{Situer chaque zéro par rapport aux rangs.}
$5{,}60=\res{5{,}6}$ : le zéro final des centièmes ne change rien. $5{,}06\ne5{,}6$ : $\res{\text{différents}}$. $07{,}2=\res{7{,}2}$ : le zéro de gauche est inutile.
\end{corrige}\end{blocLivre}
\begin{blocLivre}{cor-9-11}\begin{corrige}{11}{Le podium du saut}\demarche{Égaliser le nombre de chiffres avant de comparer.}
$3{,}5=3{,}50$ dépasse $3{,}45$, qui dépasse $3{,}09$ : $\res{\text{Maya, puis Lina, puis Inès}}$. L’ordre décroissant des longueurs est $3{,}5>3{,}45>3{,}09$.
\end{corrige}\end{blocLivre}
\begin{blocLivre}{cor-9-12}\begin{corrige}{12}{La toise de la classe}\demarche{Penser chaque taille en centimètres.}
En centièmes : $129$, $133$, $138$, $140$. L’ordre croissant est $\res{1{,}29<1{,}33<1{,}38<1{,}4}$.
\end{corrige}\end{blocLivre}
\begin{blocLivre}{cor-9-13}\begin{corrige}{13}{Toujours un intermédiaire ?}\demarche{Passer aux centièmes, voire aux rangs suivants.}
Par exemple $\res{6{,}24}$ et $\res{6{,}27}$. $\res{\text{Oui}}$ : en écrivant les deux nombres avec un rang de plus, des écritures intermédiaires apparaissent toujours ; la demi-droite n’a pas de « trous ».
\end{corrige}\end{blocLivre}
\begin{blocLivre}{cor-9-14}\begin{corrige}{14}{L’étiquette effacée}\demarche{Donner la même longueur décimale aux deux prix.}
$4{,}5=4{,}50$ : cinquante centimes contre cinq centimes. Le paquet à $\res{4{,}05}$ euros est le moins cher, avec $4{,}50-4{,}05=0{,}45$ euro d’écart. On compare des rangs, pas des écritures.
\end{corrige}\end{blocLivre}
\begin{blocLivre}{cor-9-15}\begin{corrige}{15}{L’erreur de rangement}\demarche{Comparer les dixièmes avec les dixièmes.}
Aya compare « 34 » et « 5 » comme des entiers ; or il faut comparer rang par rang : 3 dixièmes contre 5 dixièmes, donc $\res{12{,}34<12{,}5}$. Les arrondis sont $\res{12}$ et $\res{13}$ : $12{,}5$ est à égale distance et s’arrondit au-dessus.
\end{corrige}\end{blocLivre}
